% =====================================================================
%  USTH BACHELOR THESIS  --  Department of ICT  --  2026
%  Formatted to match the official "2026 ICT Bachelor thesis template".
%
%  TIP - Tap Into Prosperity: An Offline-First Personal Finance
%  Management Application for Android with Real-Time Cloud Synchronization
%
%  COMPILE ON OVERLEAF WITH:  Menu -> Compiler -> XeLaTeX  (REQUIRED)
%  -------------------------------------------------------------------
%  TODO before submitting:
%   * Replace every [PLACEHOLDER] (student name, ID, supervisors, date).
%   * Drop in the USTH logo on the title page (\includegraphics ...).
%   * Replace each \figplaceholder{...} with \includegraphics{...} and
%     upload the matching image (paths from docs/diagrams/ are noted).
%   * Body font is TeX Gyre Termes (Times New Roman-compatible, bundled
%     with Overleaf); headings use the template's navy "Heading 1" look.
% =====================================================================

\documentclass[13pt,a4paper]{extarticle}

% ---- Fonts & language (XeLaTeX) ----
\usepackage{fontspec}
\setmainfont{TeX Gyre Termes}      % Times New Roman-compatible metrics
\usepackage{polyglossia}
\setdefaultlanguage{english}

% ---- Page geometry ----
\usepackage[top=2.5cm,bottom=2.5cm,left=3.0cm,right=2.5cm,
            headsep=0.6cm,footskip=1.0cm]{geometry}

% ---- Line spacing 1.3, justified ----
\usepackage{setspace}
\setstretch{1.3}
\usepackage{microtype}

% ---- Math, tables, graphics, lists ----
\usepackage{amsmath,amssymb}
\usepackage{graphicx}
\usepackage{booktabs}
\usepackage{array}
\usepackage{longtable}
\usepackage{tabularx}
\usepackage{multirow}
\usepackage{xcolor}
\usepackage{enumitem}
\usepackage{float}

% ---- Template heading colour (Word "Heading 1" navy) ----
\definecolor{usthblue}{RGB}{31,73,125}

% ---- Captions: "Table 1: ..." / "Figure 1: ...", bold, centred ----
\usepackage{caption}
\captionsetup{labelsep=colon, font=bf, justification=centering,
              singlelinecheck=true, skip=6pt}

% ---- Headers / footers: page number bottom-right (template style) ----
\usepackage{fancyhdr}
\pagestyle{fancy}
\fancyhf{}
\fancyfoot[R]{\thepage}
\renewcommand{\headrulewidth}{0pt}
\fancypagestyle{plain}{\fancyhf{}\fancyfoot[R]{\thepage}\renewcommand{\headrulewidth}{0pt}}

% ---- Source-code listings ----
\usepackage{listings}
\definecolor{codegray}{gray}{0.5}
\definecolor{codekw}{RGB}{0,0,180}
\definecolor{codecm}{RGB}{0,128,0}
\definecolor{codestr}{RGB}{170,0,40}
\definecolor{codebg}{RGB}{248,248,248}
\lstset{
  basicstyle=\ttfamily\footnotesize,
  numbers=left, numberstyle=\tiny\color{codegray}, numbersep=8pt,
  keywordstyle=\color{codekw}\bfseries,
  commentstyle=\color{codecm}\itshape,
  stringstyle=\color{codestr},
  backgroundcolor=\color{codebg},
  frame=single, framesep=4pt, rulecolor=\color{codegray},
  breaklines=true, showstringspaces=false, tabsize=2,
  captionpos=b, columns=fullflexible, keepspaces=true
}

% ---- Sectioning: I/ , II/ ... navy heading + rule (template look) ----
\renewcommand{\thesection}{\Roman{section}/}
\renewcommand{\thesubsection}{\arabic{subsection}}
\renewcommand{\thesubsubsection}{\arabic{subsection}.\arabic{subsubsection}}
\setcounter{secnumdepth}{2}
\setcounter{tocdepth}{2}

\usepackage{titlesec}
\titleformat{\section}
  {\color{usthblue}\normalfont\fontsize{16}{19}\bfseries}
  {\thesection}{0.4em}{}
  [{\color{usthblue}\titlerule[1.4pt]}]
\titlespacing*{\section}{0pt}{16pt}{8pt}
\titleformat{\subsection}
  {\color{usthblue}\normalfont\fontsize{13.5}{16}\bfseries}
  {\thesubsection.}{0.5em}{}
\titlespacing*{\subsection}{0pt}{12pt}{4pt}
\titleformat{\subsubsection}
  {\normalfont\fontsize{13}{16}\bfseries\itshape}
  {\thesubsubsection}{0.5em}{}
\titlespacing*{\subsubsection}{0pt}{10pt}{3pt}

% ---- List-of titles & ToC name to match the template ----
\renewcommand{\contentsname}{Table of Contents}
\renewcommand{\listtablename}{LIST OF TABLES}
\renewcommand{\listfigurename}{LIST OF FIGURES}

% ---- Hyperref last ----
\usepackage[hidelinks,colorlinks=true,linkcolor=black,citecolor=blue,urlcolor=blue]{hyperref}

% ---- Helper: framed figure placeholder ----
\newcommand{\figplaceholder}[2][0.85]{%
  \fbox{\begin{minipage}[c][4.0cm][c]{#1\textwidth}
    \centering\itshape\small #2
  \end{minipage}}}

\begin{document}

% =====================================================================
%  TITLE PAGE (framed, per template)
% =====================================================================
\thispagestyle{empty}
\begin{titlepage}
\setlength{\fboxsep}{14pt}
\noindent\fbox{%
\begin{minipage}[c][\dimexpr\textheight-2\fboxsep-2\fboxrule\relax][t]{\dimexpr\textwidth-2\fboxsep-2\fboxrule\relax}
\centering
\vspace{0.4cm}
{\normalsize UNIVERSITY OF SCIENCE AND TECHNOLOGY OF HANOI\par}
{\normalsize\bfseries DEPARTMENT OF INFORMATION AND COMMUNICATION TECHNOLOGY\par}
\vspace{1.4cm}

% --- USTH logo: upload usth_logo.png next to this .tex ---
\IfFileExists{USTH_logo.png}
  {\includegraphics[width=0.45\textwidth]{USTH_logo.png}}
  {\figplaceholder[0.4]{[USTH logo here -- upload USTH\_logo.png]}}\par
\vspace{1.4cm}

{\fontsize{30}{34}\bfseries BACHELOR THESIS\par}
\vspace{0.9cm}
{\large By\par}
{\large [Student's Full Name]\par}
\vspace{1.0cm}

{\large Title:\par}
\vspace{0.2cm}
{\large\bfseries TIP -- Tap Into Prosperity: An Offline-First Personal\\
Finance Management Application for Android with\\
Real-Time Cloud Synchronization\par}
\vspace{2.6cm}

\begin{flushleft}
\hspace{1.5cm}\begin{tabular}{@{}ll@{}}
External Supervisor: & Prof. Dr. [External Supervisor Name] \\[6pt]
Internal Supervisor: & Prof. Dr. [Internal Supervisor Name] \\
\end{tabular}
\end{flushleft}
\vfill
{\large\bfseries Hanoi, [Month] -- 2026\par}
\vspace{0.4cm}
\end{minipage}}
\end{titlepage}

% =====================================================================
%  FRONT MATTER
% =====================================================================
\pagenumbering{arabic}

% ---------------------- CERTIFICATION --------------------------------
\thispagestyle{fancy}
\vspace*{1cm}
\begin{center}\textbf{To whom it may concern,}\end{center}
\vspace{0.8cm}
\noindent I, \dotfill\ (supervisor's name), certify that the thesis / internship report of
Mr/Ms.\ \dotfill\ (student's name) is qualified to be presented in the Internship Jury
2025--2026.
\vspace{2.0cm}

\begin{flushright}
Hanoi, [date]\\[4pt]
\textbf{Supervisor's signature}
\end{flushright}
\clearpage

% ---------------------- TABLE OF CONTENTS ----------------------------
\tableofcontents
\clearpage

% ---------------------- ACKNOWLEDGEMENTS -----------------------------
\section*{ACKNOWLEDGEMENTS}
\addcontentsline{toc}{section}{ACKNOWLEDGEMENTS}

I would like to express my sincere gratitude to all those who gave me the possibility to
complete this thesis, both for the specific guidance and the general academic support I received
during my studies at the University of Science and Technology of Hanoi.

Foremost, I would like to send my honest thanks to my supervisor,
\textbf{[Supervisor Full Name]}, for the invaluable guidance, patience, and constructive feedback
provided throughout this work. I am also deeply grateful to the lecturers of the Department of
Information and Communication Technology for the knowledge they have given me. Finally, I thank my
family and friends for their continuous encouragement and support.
\clearpage

% ---------------------- LIST OF ABBREVIATIONS ------------------------
\section*{LIST OF ABBREVIATIONS}
\addcontentsline{toc}{section}{LIST OF ABBREVIATIONS}
\vspace{0.4cm}
\renewcommand{\arraystretch}{1.25}
\begin{longtable}{@{\hspace{1cm}}p{2.6cm}p{9.5cm}@{}}
AES      & Advanced Encryption Standard \\
API      & Application Programming Interface \\
BCrypt   & Password-hashing function based on the Blowfish cipher \\
CRDT     & Conflict-free Replicated Data Type \\
CRUD     & Create, Read, Update, Delete \\
DAO      & Data Access Object \\
EWMA     & Exponentially Weighted Moving Average \\
FK       & Foreign Key \\
FR / NFR & Functional / Non-Functional Requirement \\
HMAC     & Hash-based Message Authentication Code \\
HTTPS    & Hypertext Transfer Protocol Secure \\
ICT      & Information and Communication Technology \\
JPA      & Jakarta Persistence API \\
JWT      & JSON Web Token \\
LLM      & Large Language Model \\
LWW      & Last-Write-Wins (conflict-resolution strategy) \\
MVVM     & Model--View--ViewModel \\
OCR      & Optical Character Recognition \\
OLS      & Ordinary Least Squares (linear regression) \\
ORM      & Object--Relational Mapping \\
REST     & Representational State Transfer \\
SDK      & Software Development Kit \\
SQL      & Structured Query Language \\
USTH     & University of Science and Technology of Hanoi \\
UUID     & Universally Unique Identifier \\
VND      & Vietnamese Dong (currency) \\
\end{longtable}
\clearpage

% ---------------------- LIST OF TABLES -------------------------------
\addcontentsline{toc}{section}{LIST OF TABLES}
\listoftables
\clearpage

% ---------------------- LIST OF FIGURES ------------------------------
\addcontentsline{toc}{section}{LIST OF FIGURES}
\listoffigures
\clearpage

% ---------------------- ABSTRACT -------------------------------------
\section*{ABSTRACT}
\addcontentsline{toc}{section}{ABSTRACT}
\textit{(No more than 250 words, in English + 6 keywords)}
\vspace{0.4cm}

This thesis presents the design and implementation of \textbf{TIP -- Tap Into Prosperity}, a
personal finance management application for Android that lets users record transactions, manage
multiple wallets, set budgets, track savings goals and debts, and obtain financial insights, with
all amounts in Vietnamese Dong. The central objective was an application that remains fully usable
without a network connection while keeping a user's data consistent across several devices. To
achieve this, the system follows an \textit{offline-first} architecture: every write is committed
to a local Room/SQLite database first and the UI reacts immediately, while a background WorkManager
job synchronizes with the cloud when connectivity is available. The Android client (Java, MVVM +
Repository) communicates over HTTPS with a stateless Spring~Boot~4 REST backend secured by JWT and
BCrypt and backed by a serverless Neon PostgreSQL database. Three problems were solved in depth:
efficient incremental synchronization through a cursor-based delta-sync protocol; multi-device
conflict resolution using Last-Write-Wins combined with server-recomputed wallet balances; and
automated receipt capture through a hybrid OCR pipeline (on-device ML~Kit text recognition plus
server-side Gemini analysis, with a local fallback). Functional verification confirmed that
authentication, transaction management, multi-account isolation, background synchronization, and
receipt extraction operate correctly, and that the backend responds within interactive latency
bounds.

\vspace{0.4cm}
\noindent\textbf{Keywords:} offline-first, Android, data synchronization, conflict resolution,
JWT authentication, hybrid OCR.
\clearpage

% =====================================================================
%  BODY  (IMRAD, per template)
% =====================================================================

% ---------------------- I / INTRODUCTION -----------------------------
\section{INTRODUCTION}

Personal financial management has become a daily necessity, yet many popular money-tracking tools
share two weaknesses: they assume a stable network connection, and they become inconsistent when
the same account is used on more than one device. Users frequently record expenses while offline
--- on public transport, while travelling, or in areas of weak coverage --- and expect the
application to respond instantly and to reconcile their data later without loss or duplication. An
application that freezes or loses data under these conditions quickly loses the user's trust.

\subsection{Background and related concepts}
The system builds on several established ideas. The \textit{offline-first} paradigm treats the
on-device database as the primary store the user interacts with and the network as an enhancement
rather than a prerequisite~\cite{android-offline}; user actions are persisted locally and reflected
immediately (an \textit{optimistic update}) while synchronization happens in the background. On
Android, this reactive pipeline is realised with Room over SQLite, whose queries return
lifecycle-aware \texttt{LiveData} observables that refresh the interface automatically when the
underlying data changes~\cite{room}. Synchronizing two replicas efficiently requires transferring
only what has changed: a \textit{delta} sync uses a \textit{cursor} (a timestamp) to fetch records
modified since the previous synchronization, paginated safely with \textit{keyset
pagination}~\cite{keyset}. When the same record is edited on two devices, a convergent
conflict-resolution policy is needed; the simplest is \textit{Last-Write-Wins} (LWW), whereas
multi-user collaboration would call for more complex Conflict-free Replicated Data
Types~\cite{crdt}. Security relies on stateless authentication with JSON Web
Tokens~\cite{jwt} and BCrypt password hashing~\cite{bcrypt}, while automated receipt entry combines
on-device OCR via ML~Kit~\cite{mlkit} with semantic interpretation by a Large Language
Model~\cite{gemini}. The overall direction of this work was also informed by existing open-source
receipt-scanning expense trackers, in particular ScanExpense~\cite{scanexpense}, which the present
project builds upon by adding an offline-first synchronization layer, multi-device conflict
resolution, and a stateless authenticated backend.

\subsection{Problem statement and motivation}
These observations motivate the central research question of this work: \emph{how can a mobile
personal-finance application provide an instantaneous, always-available user experience while still
guaranteeing data integrity across multiple devices?} Answering it requires reconciling a local
replica (optimised for responsiveness) with an authoritative cloud replica (optimised for
consistency) through a synchronization mechanism that is efficient, lossless, and resilient to
intermittent connectivity, all under strict security and per-user isolation constraints.

% ---------------------- II / OBJECTIVES ------------------------------
\section{OBJECTIVES}

The scientific objective of this work is to design and build an offline-first personal finance
application whose local database is the source of truth for the user experience and whose cloud
database is the source of truth for the system across devices; the strategy is to commit every
write locally for instant feedback and to reconcile devices through a background, cursor-based
delta-sync protocol with deterministic conflict resolution.

\noindent The concrete objectives are:
\begin{enumerate}[leftmargin=2em,itemsep=2pt]
  \item Implement a complete Android client supporting transactions (income, expense, transfer),
        multi-wallet management, budgets, goals, debts, dashboards, and spending analytics, fully
        operable offline.
  \item Design and implement a secure, stateless REST backend with JWT authentication, BCrypt
        hashing, and per-user data isolation.
  \item Engineer a bandwidth-efficient, fault-tolerant synchronization mechanism (cursor-based
        delta sync with Last-Write-Wins conflict resolution).
  \item Integrate automated receipt scanning through a hybrid OCR pipeline.
  \item Verify the system functionally and evaluate the responsiveness of its core services.
\end{enumerate}

\noindent\textbf{Scope and subjects.} The subject of study is an end-user personal finance
application together with its supporting backend. The scope covers the Android client (minimum
API~26), the Spring~Boot backend, and the Neon PostgreSQL database. A single-user-per-account model
(one user, multiple devices) is targeted; collaborative multi-user accounts, web/iOS clients, and
production deployment are treated as future work.

% ---------------------- III / MATERIALS AND METHODS ------------------
\section{MATERIALS AND METHODS}

This section presents the materials (technologies and requirements) and the methods (architecture,
data model, and the synchronization, security, and OCR designs) used to build the system. The
design proceeds top-down: \emph{requirements $\rightarrow$ use cases $\rightarrow$ architecture and
component design}.

\subsection{Materials: technologies and requirements}
The implementation technologies are summarized in Table~\ref{tab:stack}, and the principal software
and hardware requirements in Table~\ref{tab:sysreq}.

\begin{table}[H]
\centering
\caption{Implementation technology stack}
\label{tab:stack}
\renewcommand{\arraystretch}{1.2}
\begin{tabularx}{\textwidth}{@{}l X@{}}
\toprule
\textbf{Area} & \textbf{Technology} \\
\midrule
Android  & Java, Room (schema v18), Retrofit~2, OkHttp~4, Gson, LiveData/ViewModel, Navigation Component, WorkManager, Material Design~3, CameraX + ML~Kit (OCR), MPAndroidChart, EncryptedSharedPreferences. \\
Backend  & Spring Boot 4.0.5, Spring Security, Spring Data JPA, Hibernate, Lombok, jjwt 0.12.6, BCrypt, Google Gemini API. \\
Database & Neon PostgreSQL (cloud) + Room/SQLite (local cache). \\
Build    & Gradle 8.13 (AGP 8.13.2), Java~21. \\
\bottomrule
\end{tabularx}
\end{table}

\begin{table}[H]
\centering
\caption{Representative software and hardware requirements}
\label{tab:sysreq}
\renewcommand{\arraystretch}{1.2}
\begin{tabularx}{\textwidth}{@{}l X@{}}
\toprule
\textbf{Component} & \textbf{Requirement} \\
\midrule
Android client OS  & Android 7.0 (API~24) minimum; API~26+ recommended. \\
Client hardware    & 2\,GB RAM min (3\,GB+ recommended); $\sim$100\,MB storage; rear camera for scanning. \\
Backend runtime    & JDK/JRE 21; Spring Boot 4.0; 512\,MB--1\,GB RAM minimum. \\
Database           & Neon PostgreSQL (serverless cloud); Room/SQLite on device. \\
Development machine & 8\,GB RAM min (16\,GB+ recommended); virtualization for the emulator. \\
\bottomrule
\end{tabularx}
\end{table}

\subsection{Functional and non-functional requirements}
The functional requirements (Table~\ref{tab:fr}) specify what the system must do; the
non-functional requirements (Table~\ref{tab:nfr}) constrain how well it must do it.

\begin{table}[H]
\centering
\caption{Functional requirements of the TIP system}
\label{tab:fr}
\renewcommand{\arraystretch}{1.15}
\begin{tabularx}{\textwidth}{@{}l l X@{}}
\toprule
\textbf{ID} & \textbf{Group} & \textbf{Requirement} \\
\midrule
FR-01 & Auth        & Register a new account. \\
FR-02 & Auth        & Log in with email/password and receive a JWT. \\
FR-03 & Auth        & Log out; switching account wipes local data. \\
FR-04 & Auth        & Expire the session automatically on invalid/expired token. \\
FR-05 & Transaction & Record income, expense, and transfer transactions. \\
FR-06 & Transaction & Create, edit, and soft-delete transactions. \\
FR-07 & Transaction & Support recurring transactions. \\
FR-08 & Transaction & View transaction history and details. \\
FR-09 & Receipt     & Scan a receipt using CameraX + ML~Kit OCR. \\
FR-10 & Receipt     & Analyse the receipt (Gemini) and pre-fill a transaction. \\
FR-11 & Manage      & Manage wallets (CASH / BANK / E-WALLET / INVESTMENT). \\
FR-12 & Manage      & Manage income/expense categories. \\
FR-13 & Manage      & Manage budgets. \\
FR-14 & Manage      & Manage savings goals. \\
FR-15 & Manage      & Manage debts and loans. \\
FR-16 & View        & Dashboard of total assets, net worth, recent transactions. \\
FR-17 & View        & Spending analytics (bar and pie charts). \\
FR-18 & View        & Financial insights (on-device AI + Gemini enrichment). \\
FR-19 & View        & Home-screen widget (balance + three recent transactions). \\
FR-20 & Sync        & Background synchronization respecting foreign-key order. \\
FR-21 & Sync        & Cursor-based delta synchronization. \\
FR-22 & Sync        & Last-Write-Wins conflict resolution. \\
\bottomrule
\end{tabularx}
\end{table}

\begin{table}[H]
\centering
\caption{Non-functional requirements of the TIP system}
\label{tab:nfr}
\renewcommand{\arraystretch}{1.2}
\begin{tabularx}{\textwidth}{@{}l l X@{}}
\toprule
\textbf{ID} & \textbf{Attribute} & \textbf{Requirement} \\
\midrule
NFR-01 & Availability    & Offline-first: writes hit Room first; the app runs without network. \\
NFR-02 & Security        & JWT Bearer auth, BCrypt hashing, EncryptedSharedPreferences (AES256-GCM). \\
NFR-03 & Data isolation  & Four-layer multi-account isolation (schema/write/async/switch). \\
NFR-04 & Responsiveness  & Reactive UI (LiveData) with optimistic updates. \\
NFR-05 & Localization    & Vietnamese interface, currency in VND. \\
NFR-06 & Confidentiality & HTTPS for all client--server communication. \\
\bottomrule
\end{tabularx}
\end{table}

\subsection{Use case analysis}
The system has two actors: the \textbf{User} (end user) and the \textbf{Backend} (Spring Boot with
Gemini and PostgreSQL). Figure~\ref{fig:usecase} gives the overall use-case view, in which
\emph{Record Income/Expense/Transfer} inherit from an abstract \emph{Record Transaction}, the
\emph{Record} and \emph{Manage} use cases \textit{include} \emph{Synchronize Data}, and \emph{Scan
Receipt} \textit{includes} \emph{Analyze Invoice} and \textit{extends} \emph{Record Transaction}.
Table~\ref{tab:uc01} specifies the authentication use case as a representative example.

\begin{figure}[H]
\centering
\IfFileExists{MoneyTracker_UseCase.png}
  {\includegraphics[width=0.9\textwidth]{MoneyTracker_UseCase.png}}
  {\figplaceholder{Figure placeholder --- Overall use case diagram.\\[2pt]
   Upload \texttt{MoneyTracker\_UseCase.png}.}}
\caption{Overall use case diagram of the TIP system}
\label{fig:usecase}
\end{figure}

\begin{table}[H]
\centering
\caption{Use case UC-01 --- Authenticate Account}
\label{tab:uc01}
\renewcommand{\arraystretch}{1.25}
\begin{tabularx}{\textwidth}{@{}l X@{}}
\toprule
\textbf{Item} & \textbf{Description} \\
\midrule
Actors        & User, Backend. \\
Precondition  & App installed; network available for the first login. \\
Main flow     & (1) User enters email/password; (2) client calls the login endpoint; (3) backend verifies with BCrypt; (4) backend returns a JWT; (5) client stores it in EncryptedSharedPreferences; (6) navigate to the main screen. \\
Alt. flow     & No account $\rightarrow$ register; account changed $\rightarrow$ clear local tables. \\
Exception     & Wrong credentials (401) $\rightarrow$ error; later expiry $\rightarrow$ session-expired $\rightarrow$ Login. \\
Postcondition & A JWT is held in encrypted storage; the security context is ready. \\
\bottomrule
\end{tabularx}
\end{table}

\subsection{System architecture}
The system adopts a layered client--server architecture (Figure~\ref{fig:arch}). The local Room
database is the source of truth for the \emph{user experience}, whereas the cloud PostgreSQL
database is the source of truth for the \emph{system} across devices; the two are reconciled by the
background synchronization process. The Android client applies MVVM with the Repository pattern:
Activities/Fragments only render UI; \texttt{AppViewModel} (a reactive ``financial engine'' built
on \texttt{MediatorLiveData}) computes total assets, net worth, and budget consumption, all filtered
by \texttt{userId}; repositories coordinate the local Room DAOs and the remote Retrofit APIs. The
backend follows the conventional Controller--Service--Repository tiers over Spring Data JPA.

\begin{figure}[H]
\centering
\IfFileExists{architecture_layered.png}
  {\includegraphics[width=0.6\textwidth]{architecture_layered.png}}
  {\figplaceholder{Figure placeholder --- Layered architecture.\\[2pt]
   Upload \texttt{architecture\_layered.png}.}}
\caption{High-level layered architecture of the TIP system}
\label{fig:arch}
\end{figure}

\subsection{Data model and integrity}
Six core entities are stored both locally and on the server: \texttt{Transaction}, \texttt{Wallet}
(named \texttt{Account} on the server), \texttt{Category}, \texttt{Budget}, \texttt{Goal}, and
\texttt{DebtLoan}; the server adds \texttt{User} and \texttt{Session}. Every data entity carries a
\texttt{user\_id}, an \texttt{updatedAt} timestamp (for delta sync), and a soft-delete marker
(\texttt{deletedAt}/\texttt{isDeleted}). Table~\ref{tab:schema} lists the principal tables.

\begin{table}[H]
\centering
\caption{Principal database tables (Neon PostgreSQL)}
\label{tab:schema}
\renewcommand{\arraystretch}{1.2}
\begin{tabularx}{\textwidth}{@{}l X@{}}
\toprule
\textbf{Table} & \textbf{Key columns} \\
\midrule
\texttt{users}        & id (UUID), email, password\_hash (BCrypt), full\_name. \\
\texttt{accounts}     & id, user\_id (FK), name, type, balance, deleted\_at, updated\_at. \\
\texttt{categories}   & id, user\_id (FK), name, type, icon, is\_system. \\
\texttt{transactions} & id, user\_id, account\_id, category\_id, amount (CHECK $>0$), type, transaction\_date, is\_recurring, deleted\_at, updated\_at. \\
\texttt{budgets}      & id, user\_id, category\_id, amount, spent\_amount, period. \\
\texttt{goals}        & id, user\_id, name, target\_amount, current\_amount, target\_date. \\
\texttt{debts}        & id, user\_id, contact\_name, amount, type, due\_date. \\
\texttt{sessions}     & id, user\_id, refresh-token hash (SHA-256), device, timestamps. \\
\bottomrule
\end{tabularx}
\end{table}

Integrity is enforced on both sides. On the server, every balance-changing operation is atomic:
persisting a transaction and updating the wallet balance must succeed or fail together
(Listing~\ref{lst:txn}). A \texttt{CHECK(amount > 0)} constraint and foreign keys prevent invalid
and orphaned records. On the client, Room is itself an ACID database, so each write completes before
the UI reacts.

\begin{lstlisting}[language=Java,caption={Server-side atomic transaction creation (simplified).},label={lst:txn}]
@Transactional
public TransactionResponse createTransaction(CreateTransactionRequest req) {
    if (type == INCOME) account.setBalance(balance.add(amount));
    else                account.setBalance(balance.subtract(amount));
    accountRepository.save(account);                  // (1) update balance
    return toResponse(transactionRepository.save(tx)); // (2) save transaction
}                                                     // (1)+(2) atomic
\end{lstlisting}

The system uses \textit{soft delete} (flag rather than remove a row) so that a deletion becomes a
\textit{tombstone} that can be propagated through synchronization to other devices, preventing the
``zombie resurrection'' in which a hard-deleted local row is re-downloaded on the next pull, and so
that the deletion --- carrying a fresh \texttt{updatedAt} --- participates in the same Last-Write-Wins
mechanism as ordinary edits.

\subsection{Synchronization method}
Following the optimistic write path, \texttt{WorkManager} (requiring connectivity and using
exponential backoff) pushes unsynced records and then pulls deltas in foreign-key order
(\texttt{categories} $\rightarrow$ \texttt{accounts} $\rightarrow$ \texttt{budgets} $\rightarrow$
\texttt{goals} $\rightarrow$ \texttt{transactions}). Each entity exposes a \texttt{/delta} endpoint
and owns a server-issued cursor stored in \texttt{SyncPrefs}; the cursor must be server-issued
rather than read from the device clock to avoid clock-skew gaps. One delta cycle is governed by the
three parameters in Table~\ref{tab:delta}, and the server query selects the half-open window
$(\textit{since}, \textit{until}]$ with stable keyset ordering (Listing~\ref{lst:delta}).
Figure~\ref{fig:deltaseq} shows the full sequence.

\begin{table}[H]
\centering
\caption{Parameters that make delta sync lossless and duplicate-free}
\label{tab:delta}
\renewcommand{\arraystretch}{1.25}
\begin{tabularx}{\textwidth}{@{}l X@{}}
\toprule
\textbf{Parameter} & \textbf{Role} \\
\midrule
\texttt{updatedSince} (cursor) & Return only rows with \texttt{updatedAt > cursor}; the first sync uses an epoch cursor (full download). \\
\texttt{untilTimestamp} (snapshot) & The first page pins a server snapshot time reused by every later page (\texttt{updatedAt <= until}); mid-pagination edits arrive on the next cycle. \\
\texttt{lastUpdatedAt} + \texttt{lastId} & Keyset pagination on \texttt{(updatedAt, id)} instead of \texttt{OFFSET}; \texttt{id} breaks ties on equal timestamps. \\
\bottomrule
\end{tabularx}
\end{table}

\begin{lstlisting}[language=SQL,caption={Keyset delta query (simplified).},label={lst:delta}]
SELECT t FROM Transaction t WHERE t.user_id = :uid
  AND t.updatedAt > :since AND t.updatedAt <= :until      -- window (since, until]
  AND (:lastTs IS NULL OR t.updatedAt > :lastTs            -- keyset pagination
       OR (t.updatedAt = :lastTs AND t.id > :lastId))
ORDER BY t.updatedAt ASC, t.id ASC                         -- total, stable order
\end{lstlisting}

\begin{figure}[H]
\centering
\figplaceholder{Figure placeholder --- Delta-sync sequence diagram.\\[2pt]
Replace with \texttt{docs/diagrams/SEQ04\_DeltaSync\_Cursor.png}.}
\caption{Cursor-based delta synchronization sequence}
\label{fig:deltaseq}
\end{figure}

The query returns a \texttt{Slice<T>} rather than a \texttt{Page<T>}, omitting \texttt{COUNT(*)};
soft-deleted rows are intentionally included so deletions propagate. When the same record is edited
on two replicas, the version with the newer \texttt{updatedAt} wins (Last-Write-Wins). The decision
is made null-safely (Listing~\ref{lst:lww}); the edge cases handled explicitly are listed in
Table~\ref{tab:lww}.

\begin{lstlisting}[language=Java,caption={Null-safe Last-Write-Wins decision (simplified).},label={lst:lww}]
boolean clientWins = item.getClientUpdatedAt() != null
    && (existing.getUpdatedAt() == null
        || item.getClientUpdatedAt().isAfter(existing.getUpdatedAt()));
if (!clientWins) { responseList.add(toResponse(existing)); return 0; } // SERVER wins
existing.setUpdatedAt(item.getClientUpdatedAt());                      // CLIENT wins
\end{lstlisting}

\begin{table}[H]
\centering
\caption{Last-Write-Wins edge cases and decisions}
\label{tab:lww}
\renewcommand{\arraystretch}{1.2}
\begin{tabularx}{\textwidth}{@{}X l X@{}}
\toprule
\textbf{Situation} & \textbf{Decision} & \textbf{Rationale} \\
\midrule
client timestamp newer            & Client wins  & keep the user's latest change. \\
server timestamp newer            & Server wins  & client overwrites local. \\
client timestamp null (legacy)    & Server wins  & avoid legacy overwrite of good data. \\
server timestamp null (old row)   & Client wins  & the old row was never stamped. \\
deleted but client resends        & Re-acknowledge & avoid an infinite retry loop. \\
duplicate new (legacy, id null)   & Return existing & prevent double-insert on resend. \\
\bottomrule
\end{tabularx}
\end{table}

Because Last-Write-Wins on a running balance counter would drift under concurrent edits, the system
never mutates a counter; it \emph{recomputes} each wallet balance from the sum of its transactions
(Equation~\ref{eq:balance}), and derives net worth from the included wallets less debts plus loans
(Equation~\ref{eq:networth}):

\begin{equation}
\label{eq:balance}
\textit{balance} \;=\; \sum_{t \in \text{income}} t.\textit{amount} \;-\!\! \sum_{t \in \text{expense} \cup \text{transfer}} t.\textit{amount}
\end{equation}
\begin{equation}
\label{eq:networth}
\textit{netWorth} \;=\; \textit{totalAssets} \;-\; \textit{debts} \;+\; \textit{loans}
\end{equation}

\subsection{Security and data isolation}
Authentication is stateless and based on JWT (Figure~\ref{fig:loginseq}): at login the backend
verifies the password with BCrypt and issues a token signed with HMAC-SHA; thereafter an OkHttp
interceptor attaches the token to every request and a filter verifies it. The token is stored in
\texttt{EncryptedSharedPreferences} (AES256-GCM, Android Keystore). For multi-device use, each login
also issues a rotating refresh token (stored only as a SHA-256 hash) tracked in the
\texttt{sessions} table, enabling immediate remote revocation. User data is isolated across four
layers (Table~\ref{tab:isolation}).

\begin{figure}[H]
\centering
\IfFileExists{SEQ05_Login_JWT.png}
  {\includegraphics[width=0.85\textwidth]{SEQ05_Login_JWT.png}}
  {\figplaceholder{Figure placeholder --- Login/JWT sequence diagram.\\[2pt]
   Upload \texttt{SEQ05\_Login\_JWT.png}.}}
\caption{Stateless JWT login sequence}
\label{fig:loginseq}
\end{figure}

\begin{table}[H]
\centering
\caption{Four-layer multi-account data isolation}
\label{tab:isolation}
\renewcommand{\arraystretch}{1.2}
\begin{tabularx}{\textwidth}{@{}l X@{}}
\toprule
\textbf{Layer} & \textbf{Mechanism} \\
\midrule
Schema        & a \texttt{user\_id} column on every entity; all DAO queries filter by it. \\
Write         & \texttt{userId} is stamped before every insert (user action and sync). \\
Async guard   & in-flight responses are discarded if \texttt{userId} changed since the request. \\
Account switch & cancel HTTP requests $\rightarrow$ cancel worker $\rightarrow$ wipe Room $\rightarrow$ reset cursors. \\
\bottomrule
\end{tabularx}
\end{table}

\subsection{Hybrid OCR method}
Receipt extraction uses a hybrid pipeline (Figure~\ref{fig:ocrseq}): (1) CameraX captures the image
and ML~Kit recognises text \emph{on-device}; (2) only the extracted text --- never the image --- is
posted to the backend with the JWT; (3) \texttt{GeminiService} applies prompt engineering and a
structured-output schema (low \texttt{temperature}) to coerce well-formed JSON, prioritising the
grand total and rejecting noise such as ``cash tendered''; (4) the fields \texttt{\{amount,
shopName, date, note\}} pre-fill the transaction form. This keeps the receipt image private,
offloads image processing to the device, and improves semantic accuracy over rigid regular
expressions; if the network or API fails, the client falls back to a local parser.

\begin{figure}[H]
\centering
\IfFileExists{SEQ02_ScanReceipt_HybridOCR.png}
  {\includegraphics[width=0.82\textwidth]{SEQ02_ScanReceipt_HybridOCR.png}}
  {\figplaceholder{Figure placeholder --- Hybrid OCR sequence diagram.\\[2pt]
   Upload \texttt{SEQ02\_ScanReceipt\_HybridOCR.png}.}}
\caption{Hybrid OCR receipt-analysis pipeline}
\label{fig:ocrseq}
\end{figure}

\subsection{On-device AI insight method}
Before optionally enriching with Gemini, statistical analysis runs entirely on the device
(Table~\ref{tab:insight}). For anomaly detection, a transaction amount $x$ is flagged when its
Z-score exceeds a threshold $\tau$ (Equation~\ref{eq:zscore}), where $\mu$ and $\sigma$ are the mean
and standard deviation of recent spending:

\begin{equation}
\label{eq:zscore}
z \;=\; \frac{x - \mu}{\sigma}, \qquad \text{flag if } |z| > \tau .
\end{equation}

\begin{table}[H]
\centering
\caption{On-device AI insight modules}
\label{tab:insight}
\renewcommand{\arraystretch}{1.2}
\begin{tabularx}{\textwidth}{@{}l X@{}}
\toprule
\textbf{Module} & \textbf{Method} \\
\midrule
\texttt{BudgetForecaster} & Ordinary Least Squares regression --- spending forecast. \\
\texttt{AnomalyDetector}  & Z-score --- unusual-transaction detection. \\
\texttt{PatternAnalyzer}  & EWMA --- day-of-week spending trends. \\
\texttt{GoalAdvisor}      & savings-velocity recommendations. \\
\texttt{DebtLoanAdvisor}  & upcoming due-date reminders. \\
\bottomrule
\end{tabularx}
\end{table}

% ---------------------- IV / RESULTS AND DISCUSSION ------------------
\section{RESULTS AND DISCUSSION}

\subsection{Results}
The complete system was implemented as designed and delivers the full feature set of
Table~\ref{tab:fr}. The end-to-end transaction flow behaves exactly as intended: the form delegates
to \texttt{AppViewModel.addTransaction()}, which stamps \texttt{userId}, sets \texttt{isSynced =
false} and \texttt{updatedAtMs}, and inserts into Room; Room's \texttt{LiveData} drives the
optimistic UI update before \texttt{TransactionSyncWorker} performs the push and delta-pull. Every
asynchronous sync callback is protected by a ``stale-response guard'' (Listing~\ref{lst:guard}) that
discards a response if the account changed while it was in flight. Representative screens are shown
in Figures~\ref{fig:dashboard}--\ref{fig:scan}.

\begin{lstlisting}[language=Java,caption={Stale-response guard preventing cross-account data writes.},label={lst:guard}]
final String requestUserId = tokenManager.getUserId(); // captured at send time
api.getX().enqueue(resp -> {
    if (!requestUserId.equals(tokenManager.getUserId())) return; // account changed -> drop
    // ... write to DB stamped with requestUserId ...
});
\end{lstlisting}

\begin{figure}[H]
\centering
\figplaceholder{Figure placeholder --- Dashboard screen (total assets, net worth, recent transactions).\\[2pt]
Insert the app screenshot here.}
\caption{Dashboard screen of the TIP application}
\label{fig:dashboard}
\end{figure}

\begin{figure}[H]
\centering
\figplaceholder{Figure placeholder --- Spending analytics (bar/pie charts).\\[2pt]
Insert the app screenshot here.}
\caption{Spending analytics screen}
\label{fig:analytics}
\end{figure}

\begin{figure}[H]
\centering
\figplaceholder{Figure placeholder --- Receipt scanning and extraction screen.\\[2pt]
Insert the app screenshot here.}
\caption{Receipt scanning and extraction screen}
\label{fig:scan}
\end{figure}

Functional verification (Table~\ref{tab:verif}) confirmed correct behaviour across authentication,
transaction management, delta synchronization, multi-account isolation, and receipt analysis. The
offline-first behaviour was validated by performing create/edit/delete operations with connectivity
disabled and observing that the UI stayed responsive and that all changes synchronized correctly
once connectivity returned. Backend response times, measured by repeated requests under local
testing, were within interactive bounds (tens of milliseconds for reads), with the delta protocol
transferring substantially less data per sync than a full pull.

\begin{table}[H]
\centering
\caption{Backend functional verification results}
\label{tab:verif}
\renewcommand{\arraystretch}{1.2}
\begin{tabularx}{\textwidth}{@{}l X@{}}
\toprule
\textbf{Component} & \textbf{Status} \\
\midrule
Database connection            & Connected and reachable (Neon PostgreSQL). \\
Authentication (login/register) & JWT issued and validated; BCrypt verification correct. \\
Transaction CRUD + sync        & Create/edit/soft-delete and batch sync verified. \\
Delta endpoints                & Incremental pulls return only changed rows; tombstones propagated. \\
Multi-account isolation        & Switching account wipes local data; no cross-user leakage. \\
Receipt analysis               & ML~Kit OCR + Gemini extraction returns structured fields. \\
\bottomrule
\end{tabularx}
\end{table}

\subsection{Discussion}
The results validate the central design hypothesis: committing writes locally and reconciling in the
background yields an interface that feels instantaneous yet converges to a consistent global state.
The choice of Last-Write-Wins is appropriate for this single-user-per-account setting, where genuine
conflicts are rare and arise only across one user's devices; it converges deterministically without
distributed locking, whereas CRDTs~\cite{crdt} would add considerable complexity for no practical
benefit at this scale. Deriving wallet balances from the transaction log rather than maintaining a
counter proved essential, since a counter drifts under concurrent edits while a derived balance
remains correct regardless of sync order. The hybrid OCR design likewise paid off: keeping image
processing on the device protects privacy and reduces server load, while delegating semantic
interpretation to a constrained LLM is markedly more robust than regular expressions, and the local
fallback ensures the feature never fails outright. The principal limitation observed is that the
optimistic-concurrency (409) capability exists in the backend but is not yet wired into the Android
edit screen, which currently resolves edits through the Last-Write-Wins path; the system is also
single-platform and uses a development tunnel rather than production infrastructure.

% ---------------------- V / CONCLUSION & PERSPECTIVE -----------------
\section{CONCLUSION \& PERSPECTIVE}

This thesis designed and implemented TIP -- Tap Into Prosperity, a complete offline-first personal
finance application for Android with a secure cloud backend. The system delivers transactions,
multi-wallet management, budgets, goals, debts, dashboards, analytics, insights, and a home-screen
widget, and remains fully usable without a network connection. Its main contributions are an
offline-first architecture with optimistic UI updates; a bandwidth-efficient, lossless cursor-based
delta-sync protocol using a server-pinned snapshot window and keyset pagination; a Last-Write-Wins
conflict-resolution scheme paired with server-recomputed balances that prevents value drift; a
hybrid OCR pipeline that keeps receipt images on the device while using a constrained LLM for robust
extraction; and a four-layer data-isolation model with stateless JWT authentication.

The work opens several perspectives. The optimistic-concurrency endpoint could be connected to the
edit screen; web and iOS clients could share the same backend; collaborative multi-user accounts
could be supported, which would motivate a move from Last-Write-Wins toward CRDT-based merging; the
deployment could be hardened with a managed reverse proxy and TLS termination; and the insight
engine could be extended with richer models. A periodic purge job could also reclaim space from
fully synchronized tombstones. Together these directions would evolve the present prototype into a
production-grade personal finance platform.

% =====================================================================
%  REFERENCES  (Elsevier numbered style, in order of appearance)
% =====================================================================
\section*{REFERENCES}
\addcontentsline{toc}{section}{REFERENCES}
\begin{thebibliography}{99}

\bibitem{android-offline} Google, Build an offline-first app, Android Developers Architecture Guide, 2023. \url{https://developer.android.com/topic/architecture/data-layer/offline-first} (accessed 2026).

\bibitem{room} Google, Save data in a local database using Room, Android Jetpack Documentation, 2024. \url{https://developer.android.com/training/data-storage/room} (accessed 2026).

\bibitem{keyset} M. Winand, Use the Index, Luke! -- Pagination done right (the seek method), 2017. \url{https://use-the-index-luke.com/no-offset} (accessed 2026).

\bibitem{crdt} M. Shapiro, N. Pregui\c{c}a, C. Baquero, M. Zawirski, Conflict-free Replicated Data Types, in: Proc. 13th Int. Symp. on Stabilization, Safety, and Security of Distributed Systems (SSS), Springer, 2011, pp. 386--400. \url{https://doi.org/10.1007/978-3-642-24550-3_29}

\bibitem{jwt} M. Jones, J. Bradley, N. Sakimura, JSON Web Token (JWT), RFC 7519, Internet Engineering Task Force (IETF), 2015. \url{https://doi.org/10.17487/RFC7519}

\bibitem{bcrypt} N. Provos, D. Mazi\`eres, A future-adaptable password scheme, in: Proc. USENIX Annual Technical Conference, 1999, pp. 81--91.

\bibitem{mlkit} Google, ML Kit Text Recognition v2 (on-device API), Google Developers Documentation, 2024. \url{https://developers.google.com/ml-kit/vision/text-recognition} (accessed 2026).

\bibitem{gemini} Google, Gemini API: structured output and generation configuration, Google AI for Developers, 2024. \url{https://ai.google.dev/} (accessed 2026).

\bibitem{springboot} VMware, Spring Boot Reference Documentation, version 4.0, 2024. \url{https://spring.io/projects/spring-boot} (accessed 2026).

\bibitem{scanexpense} V. Cheepurupalli, ScanExpense: a receipt-scanning expense tracker [source code], GitHub, n.d. \url{https://github.com/VignanCheepurupalli/ScanExpense} (accessed 2026).

\end{thebibliography}

% =====================================================================
%  APPENDICES
% =====================================================================
\clearpage
\section*{APPENDICES}
\addcontentsline{toc}{section}{APPENDICES}

\begin{center}\textbf{APPENDIX 1 --- Project repository and supplementary listings}\end{center}
The complete source code is organised under the package \texttt{vn.edu.usth.tip}, split into the
Android module (\texttt{app/}) and the Spring Boot module (\texttt{backend/}).
[Insert the project's own repository URL or attached media reference here.] A representative listing
of the soft-delete (tombstone) logic is shown below.

\begin{lstlisting}[language=Java,caption={Soft delete marks the row and stamps a new timestamp so LWW still applies.}]
OffsetDateTime now = OffsetDateTime.now();
tx.setDeletedAt(now);
tx.setUpdatedAt(now);     // keep LWW applicable to the deletion
tx.setClientSync(true);
transactionRepository.save(tx);
\end{lstlisting}

\vspace{0.6cm}
\begin{center}\textbf{APPENDIX 2 --- Additional diagrams}\end{center}
[Insert the database schema diagram and additional sequence diagrams from \texttt{docs/diagrams/}
here, each with a numbered caption.]

\end{document}
